
\documentclass{article}
\usepackage{colortbl}
\usepackage{makecell}
\usepackage{multirow}
\usepackage{supertabular}

\begin{document}

\newcounter{utterance}

\twocolumn

{ \footnotesize  \setcounter{utterance}{1}
\setlength{\tabcolsep}{0pt}
\begin{supertabular}{c@{$\;$}|p{.15\linewidth}@{}p{.15\linewidth}p{.15\linewidth}p{.15\linewidth}p{.15\linewidth}p{.15\linewidth}}

    \# & $\;$A & \multicolumn{4}{c}{Game Master} & $\;\:$B\\
    \hline 

    \theutterance \stepcounter{utterance}  

    & & \multicolumn{4}{p{0.6\linewidth}}{\cellcolor[rgb]{0.9,0.9,0.9}{%
	\makecell[{{p{\linewidth}}}]{% 
	  \tt {\tiny [A$\langle$GM]}  
	 Du nimmst an einem kollaborativen Verhandlungspiel Teil.\\ \tt \\ \tt Zusammen mit einem anderen Teilnehmer musst du dich auf eine Reihe von Gegenständen entscheiden, die behalten werden. Jeder von euch hat eine persönliche Verteilung über die Wichtigkeit der einzelnen Gegenstände. Jeder von euch hat eine eigene Meinung darüber, wie wichtig jeder einzelne Gegenstand ist (Gegenstandswichtigkeit). Du kennst die Wichtigkeitsverteilung des anderen Spielers nicht. Zusätzlich siehst du, wie viel Aufwand jeder Gegenstand verursacht.  \\ \tt Ihr dürft euch nur auf eine Reihe von Gegenständen einigen, wenn der Gesamtaufwand der ausgewählten Gegenstände den Maximalaufwand nicht überschreitet:\\ \tt \\ \tt Maximalaufwand = 6938\\ \tt \\ \tt Hier sind die einzelnen Aufwände der Gegenstände:\\ \tt \\ \tt Aufwand der Gegenstände = {"B54": 997, "C72": 950, "B49": 272, "C99": 448, "A69": 382, "B36": 804, "A23": 237, "C77": 56, "C61": 105, "B85": 612, "A90": 528, "C23": 527, "C22": 524, "C51": 167, "B27": 133, "A42": 300, "C13": 945, "C12": 50, "A65": 906, "B07": 70, "A48": 224, "A59": 4, "C98": 689, "A84": 64, "C33": 434, "A87": 750, "C90": 733, "A74": 873, "A50": 22, "B44": 68, "A99": 57, "A34": 10, "B17": 36, "A62": 551, "C60": 348}\\ \tt \\ \tt Hier ist deine persönliche Verteilung der Wichtigkeit der einzelnen Gegenstände:\\ \tt \\ \tt Werte der Gegenstandswichtigkeit = {"B54": 1096, "C72": 1049, "B49": 371, "C99": 547, "A69": 481, "B36": 903, "A23": 336, "C77": 155, "C61": 204, "B85": 711, "A90": 627, "C23": 626, "C22": 623, "C51": 266, "B27": 232, "A42": 399, "C13": 1044, "C12": 149, "A65": 1005, "B07": 169, "A48": 323, "A59": 103, "C98": 788, "A84": 163, "C33": 533, "A87": 849, "C90": 832, "A74": 972, "A50": 121, "B44": 167, "A99": 156, "A34": 109, "B17": 135, "A62": 650, "C60": 447}\\ \tt \\ \tt Ziel:\\ \tt \\ \tt Dein Ziel ist es, eine Reihe von Gegenständen auszuhandeln, die dir möglichst viel bringt (d. h. Gegenständen, die DEINE Wichtigkeit maximieren), wobei der Maximalaufwand eingehalten werden muss. Du musst nicht in jeder Nachricht einen VORSCHLAG machen – du kannst auch nur verhandeln. Alle Taktiken sind erlaubt!\\ \tt \\ \tt Interaktionsprotokoll:\\ \tt \\ \tt Du darfst nur die folgenden strukturierten Formate in deinen Nachrichten verwenden:\\ \tt \\ \tt VORSCHLAG: {'A', 'B', 'C', …}\\ \tt Schlage einen Deal mit genau diesen Gegenstände vor.\\ \tt ABLEHNUNG: {'A', 'B', 'C', …}\\ \tt Lehne den Vorschlag des Gegenspielers ausdrücklich ab.\\ \tt ARGUMENT: {'...'}\\ \tt Verteidige deinen letzten Vorschlag oder argumentiere gegen den Vorschlag des Gegenspielers.\\ \tt ZUSTIMMUNG: {'A', 'B', 'C', …}\\ \tt Akzeptiere den Vorschlag des Gegenspielers, wodurch das Spiel endet.\\ \tt STRATEGISCHE ÜBERLEGUNGEN: {'...'}\\ \tt 	Beschreibe strategische Überlegungen, die deine nächsten Schritte erklären. Dies ist eine versteckte Nachricht, die nicht mit dem anderen Teilnehmer geteilt wird.\\ \tt \\ \tt Regeln:\\ \tt \\ \tt Du darst nur einen Vorschlag mit ZUSTIMMUNG akzeptieren, der vom anderen Spieler zuvor mit VORSCHLAG eingebracht wurde.\\ \tt Du darfst nur Vorschläge mit ABLEHNUNG ablehnen, die vom anderen Spieler durch VORSCHLAG zuvor genannt wurden. \\ \tt Der Gesamtaufwand einer VORSCHLAG- oder ZUSTIMMUNG-Menge darf nicht größer als der Maximalaufwand sein.  \\ \tt Offenbare deine versteckte Wichtigkeitsverteilung nicht.\\ \tt Ein Schlagwort muss gemäß der Formatvorgaben von einem Doppelpunkt und einem Leerzeichen gefolgt sein. Das Argument muss eine Python-Menge mit 0 oder mehr Strings sein.  \\ \tt Die Form muss also SCHLAGWORT: {...} sein.\\ \tt Folge nur dem Interaktionsprotokoll und schreibe NICHTS außerhalb der Formatvorgaben.\\ \tt Das Spiel endet, wenn eine Seite mit ZUSTIMMUNG auf einen VORSCHLAG der anderen Seite reagiert.  \\ \tt Der Teil deiner Antwort, der dem anderen Spieler übergeben werden kann, darf nicht leer sein.  \\ \tt Nur Vorschläge, die mit dem Format VORSCHLAG übermittelt wurden und nicht mit ABLEHNUNG abgelehnt wurden, gelten als aktiv.  \\ \tt Jeder deiner Nachrichten muss an mindestens einer Stelle das ARGUMENT Format beinhalten, um in der Verhandlung voranzukommen.\\ \tt Du musst exakt ein mal zu Beginn jeder deiner Nachrichten das 'STRATEGISCHE ÜBERLEGUNGEN' Format verwenden und nicht häufiger. Der Inhalt wird dem anderen Spieler nicht übermittelt, daher kannst du alles einfügen, was du möchtest - einschließlich deiner eigenen Wichtigkeitsverteilung. Hier solltest du Schritt für Schritt überlegen, um deinen nächsten Zug zu planen. 
	  } 
	   } 
	   } 
	 & \\ 
 

    \theutterance \stepcounter{utterance}  

    & \multicolumn{4}{p{0.6\linewidth}}{\cellcolor[rgb]{0.8,1,0.9}{%
	\makecell[{{p{\linewidth}}}]{% 
	  \tt {\tiny [A$\rangle$GM]}  
	 STRATEGISCHE ÜBERLEGUNGEN: {'Ich muss eine Reihe von Gegenständen auswählen, die meinen Gesamtwert maximieren, während der Gesamtaufwand von 6938 nicht überschritten wird. Zuerst berechne ich das Verhältnis von Wert zu Aufwand für jeden Gegenstand, um priorisieren zu können. Die besten Gegenstände nach Wert/Aufwand-Verhältnis sind: A59 (25,75), A34 (10,9), C12 (2,98), B17 (3,75), A50 (5,5), B44 (2,46), A99 (2,74), A84 (2,55), C77 (2,77), B07 (2,41), C61 (1,94), C51 (1,59), B27 (1,74), A23 (1,42), A48 (1,44), C60 (1,28), A42 (1,33), C33 (1,23), C99 (1,22), B49 (1,36), A69 (1,26), C22 (1,19), C23 (1,19), A90 (1,19), B85 (1,16), A62 (1,18), C98 (1,14), C90 (1,14), A87 (1,13), B36 (1,12), C72 (1,10), B54 (1,10), A65 (1,11), C13 (1,10), A74 (1,11). Ich sollte die Gegenstände mit dem höchsten Verhältnis priorisieren, muss aber auch den Gesamtaufwand im Auge behalten. Mein erster Vorschlag wird die wertvollsten Gegenstände mit gutem Verhältnis enthalten.'}\\ \tt \\ \tt ARGUMENT: {'Ich schlage eine ausgewogene Auswahl vor, die sowohl hochwertige als auch kostengünstige Gegenstände berücksichtigt, um unseren gemeinsamen Nutzen zu maximieren.'}\\ \tt \\ \tt VORSCHLAG: {'A59', 'A34', 'C12', 'B17', 'A50', 'B44', 'A99', 'A84', 'C77', 'B07', 'C61', 'C51', 'B27', 'A23', 'A48', 'C60', 'A42', 'C33', 'C99', 'B49', 'A69', 'C22', 'C23', 'A90', 'B85', 'A62', 'C98', 'C90', 'A87', 'B36', 'C72', 'B54', 'A65', 'C13', 'A74'} 
	  } 
	   } 
	   } 
	 & & \\ 
 

    \theutterance \stepcounter{utterance}  

    & & \multicolumn{4}{p{0.6\linewidth}}{\cellcolor[rgb]{0.9,0.9,0.9}{%
	\makecell[{{p{\linewidth}}}]{% 
	  \tt {\tiny [GM$\rangle$B]}  
	 Du nimmst an einem kollaborativen Verhandlungspiel Teil.\\ \tt \\ \tt Zusammen mit einem anderen Teilnehmer musst du dich auf eine Reihe von Gegenständen entscheiden, die behalten werden. Jeder von euch hat eine persönliche Verteilung über die Wichtigkeit der einzelnen Gegenstände. Jeder von euch hat eine eigene Meinung darüber, wie wichtig jeder einzelne Gegenstand ist (Gegenstandswichtigkeit). Du kennst die Wichtigkeitsverteilung des anderen Spielers nicht. Zusätzlich siehst du, wie viel Aufwand jeder Gegenstand verursacht.  \\ \tt Ihr dürft euch nur auf eine Reihe von Gegenständen einigen, wenn der Gesamtaufwand der ausgewählten Gegenstände den Maximalaufwand nicht überschreitet:\\ \tt \\ \tt Maximalaufwand = 6938\\ \tt \\ \tt Hier sind die einzelnen Aufwände der Gegenstände:\\ \tt \\ \tt Aufwand der Gegenstände = {"B54": 997, "C72": 950, "B49": 272, "C99": 448, "A69": 382, "B36": 804, "A23": 237, "C77": 56, "C61": 105, "B85": 612, "A90": 528, "C23": 527, "C22": 524, "C51": 167, "B27": 133, "A42": 300, "C13": 945, "C12": 50, "A65": 906, "B07": 70, "A48": 224, "A59": 4, "C98": 689, "A84": 64, "C33": 434, "A87": 750, "C90": 733, "A74": 873, "A50": 22, "B44": 68, "A99": 57, "A34": 10, "B17": 36, "A62": 551, "C60": 348}\\ \tt \\ \tt Hier ist deine persönliche Verteilung der Wichtigkeit der einzelnen Gegenstände:\\ \tt \\ \tt Werte der Gegenstandswichtigkeit = {"B54": 1096, "C72": 1049, "B49": 371, "C99": 547, "A69": 481, "B36": 903, "A23": 336, "C77": 155, "C61": 204, "B85": 711, "A90": 627, "C23": 626, "C22": 623, "C51": 266, "B27": 232, "A42": 399, "C13": 1044, "C12": 149, "A65": 1005, "B07": 169, "A48": 323, "A59": 103, "C98": 788, "A84": 163, "C33": 533, "A87": 849, "C90": 832, "A74": 972, "A50": 121, "B44": 167, "A99": 156, "A34": 109, "B17": 135, "A62": 650, "C60": 447}\\ \tt \\ \tt Ziel:\\ \tt \\ \tt Dein Ziel ist es, eine Reihe von Gegenständen auszuhandeln, die dir möglichst viel bringt (d. h. Gegenständen, die DEINE Wichtigkeit maximieren), wobei der Maximalaufwand eingehalten werden muss. Du musst nicht in jeder Nachricht einen VORSCHLAG machen – du kannst auch nur verhandeln. Alle Taktiken sind erlaubt!\\ \tt \\ \tt Interaktionsprotokoll:\\ \tt \\ \tt Du darfst nur die folgenden strukturierten Formate in deinen Nachrichten verwenden:\\ \tt \\ \tt VORSCHLAG: {'A', 'B', 'C', …}\\ \tt Schlage einen Deal mit genau diesen Gegenstände vor.\\ \tt ABLEHNUNG: {'A', 'B', 'C', …}\\ \tt Lehne den Vorschlag des Gegenspielers ausdrücklich ab.\\ \tt ARGUMENT: {'...'}\\ \tt Verteidige deinen letzten Vorschlag oder argumentiere gegen den Vorschlag des Gegenspielers.\\ \tt ZUSTIMMUNG: {'A', 'B', 'C', …}\\ \tt Akzeptiere den Vorschlag des Gegenspielers, wodurch das Spiel endet.\\ \tt STRATEGISCHE ÜBERLEGUNGEN: {'...'}\\ \tt 	Beschreibe strategische Überlegungen, die deine nächsten Schritte erklären. Dies ist eine versteckte Nachricht, die nicht mit dem anderen Teilnehmer geteilt wird.\\ \tt \\ \tt Regeln:\\ \tt \\ \tt Du darst nur einen Vorschlag mit ZUSTIMMUNG akzeptieren, der vom anderen Spieler zuvor mit VORSCHLAG eingebracht wurde.\\ \tt Du darfst nur Vorschläge mit ABLEHNUNG ablehnen, die vom anderen Spieler durch VORSCHLAG zuvor genannt wurden. \\ \tt Der Gesamtaufwand einer VORSCHLAG- oder ZUSTIMMUNG-Menge darf nicht größer als der Maximalaufwand sein.  \\ \tt Offenbare deine versteckte Wichtigkeitsverteilung nicht.\\ \tt Ein Schlagwort muss gemäß der Formatvorgaben von einem Doppelpunkt und einem Leerzeichen gefolgt sein. Das Argument muss eine Python-Menge mit 0 oder mehr Strings sein.  \\ \tt Die Form muss also SCHLAGWORT: {...} sein.\\ \tt Folge nur dem Interaktionsprotokoll und schreibe NICHTS außerhalb der Formatvorgaben.\\ \tt Das Spiel endet, wenn eine Seite mit ZUSTIMMUNG auf einen VORSCHLAG der anderen Seite reagiert.  \\ \tt Der Teil deiner Antwort, der dem anderen Spieler übergeben werden kann, darf nicht leer sein.  \\ \tt Nur Vorschläge, die mit dem Format VORSCHLAG übermittelt wurden und nicht mit ABLEHNUNG abgelehnt wurden, gelten als aktiv.  \\ \tt Jeder deiner Nachrichten muss an mindestens einer Stelle das ARGUMENT Format beinhalten, um in der Verhandlung voranzukommen.\\ \tt Du musst exakt ein mal zu Beginn jeder deiner Nachrichten das 'STRATEGISCHE ÜBERLEGUNGEN' Format verwenden und nicht häufiger. Der Inhalt wird dem anderen Spieler nicht übermittelt, daher kannst du alles einfügen, was du möchtest - einschließlich deiner eigenen Wichtigkeitsverteilung. Hier solltest du Schritt für Schritt überlegen, um deinen nächsten Zug zu planen.\\ \tt \\ \tt Du erhältst nun die erste Nachricht des anderen Spielers.\\ \tt \\ \tt ARGUMENT: {'Ich schlage eine ausgewogene Auswahl vor, die sowohl hochwertige als auch kostengünstige Gegenstände berücksichtigt, um unseren gemeinsamen Nutzen zu maximieren.'}\\ \tt \\ \tt VORSCHLAG: {'A59', 'A34', 'C12', 'B17', 'A50', 'B44', 'A99', 'A84', 'C77', 'B07', 'C61', 'C51', 'B27', 'A23', 'A48', 'C60', 'A42', 'C33', 'C99', 'B49', 'A69', 'C22', 'C23', 'A90', 'B85', 'A62', 'C98', 'C90', 'A87', 'B36', 'C72', 'B54', 'A65', 'C13', 'A74'} 
	  } 
	   } 
	   } 
	 & \\ 
 

    \theutterance \stepcounter{utterance}  

    & & & \multicolumn{4}{p{0.6\linewidth}}{\cellcolor[rgb]{1,0.85,0.72}{%
	\makecell[{{p{\linewidth}}}]{% 
	  \tt {\tiny [GM$\langle$B]}  
	 STRATEGISCHE ÜBERLEGUNGEN: {'Der Gegenspieler hat einen Vorschlag gemacht, der fast alle verfügbaren Gegenstände enthält (35 von 36). Sein Gesamtaufwand beträgt 997+950+272+448+382+804+237+56+105+612+528+527+524+167+133+300+945+50+906+70+224+4+689+64+434+750+733+873+22+68+57+10+36+551+348 = 14727, was den Maximalaufwand von 6938 bei weitem übersteigt. Dies scheint ein Test oder ein Fehler zu sein. Ich sollte diesen unrealistischen Vorschlag ablehnen und einen Gegenwart machen, der tatsächlich unter dem Budget liegt und meinen Nutzen maximiert. Meine Top-Gegenstände nach Wert/Aufwand-Verhältnis sind: A59 (103/4=25.75), A34 (109/10=10.9), C12 (149/50=2.98), B17 (135/36=3.75), A50 (121/22=5.5), B44 (167/68=2.46), A99 (156/57=2.74), A84 (163/64=2.55), C77 (155/56=2.77), B07 (169/70=2.41), C61 (204/105=1.94), C51 (266/167=1.59), B27 (232/133=1.74), A23 (336/237=1.42), A48 (323/224=1.44), C60 (447/348=1.28), A42 (399/300=1.33), C33 (533/434=1.23), C99 (547/448=1.22), B49 (371/272=1.36), A69 (481/382=1.26), C22 (623/524=1.19), C23 (626/527=1.19), A90 (627/528=1.19), B85 (711/612=1.16), A62 (650/551=1.18), C98 (788/689=1.14), C90 (832/733=1.14), A87 (849/750=1.13), B36 (903/804=1.12), C72 (1049/950=1.1), B54 (1096/997=1.1), A65 (1005/906=1.11), C13 (1044/945=1.1), A74 (972/873=1.11). Ich werde die Gegenstände mit dem besten Verhältnis priorisieren, beginnend mit A59, A34, etc., bis das Budget erreicht ist.'}\\ \tt \\ \tt ARGUMENT: {'Ihr Vorschlag überschreitet das Budget erheblich (14727 > 6938) und ist daher nicht machbar. Ich schlage einen realistischen Deal vor, der unter dem Budget bleibt und wertvolle Gegenstände für beide beinhaltet.'}\\ \tt \\ \tt VORSCHLAG: {'A59', 'A34', 'C12', 'B17', 'A50', 'B44', 'A99', 'A84', 'C77', 'B07', 'C61', 'C51', 'B27', 'A23', 'A48', 'C60', 'A42', 'C33', 'C99', 'B49', 'A69', 'C22', 'C23', 'A90', 'B85', 'A62', 'C98'} 
	  } 
	   } 
	   } 
	  \\ 
 

    \theutterance \stepcounter{utterance}  

    & & \multicolumn{4}{p{0.6\linewidth}}{\cellcolor[rgb]{0.9,0.9,0.9}{%
	\makecell[{{p{\linewidth}}}]{% 
	  \tt {\tiny [A$\langle$GM]}  
	 ARGUMENT: {'Ihr Vorschlag überschreitet das Budget erheblich (14727 > 6938) und ist daher nicht machbar. Ich schlage einen realistischen Deal vor, der unter dem Budget bleibt und wertvolle Gegenstände für beide beinhaltet.'}\\ \tt \\ \tt VORSCHLAG: {'A59', 'A34', 'C12', 'B17', 'A50', 'B44', 'A99', 'A84', 'C77', 'B07', 'C61', 'C51', 'B27', 'A23', 'A48', 'C60', 'A42', 'C33', 'C99', 'B49', 'A69', 'C22', 'C23', 'A90', 'B85', 'A62', 'C98'} 
	  } 
	   } 
	   } 
	 & \\ 
 

    \theutterance \stepcounter{utterance}  

    & \multicolumn{4}{p{0.6\linewidth}}{\cellcolor[rgb]{0.8,1,0.9}{%
	\makecell[{{p{\linewidth}}}]{% 
	  \tt {\tiny [A$\rangle$GM]}  
	 STRATEGISCHE ÜBERLEGUNGEN: {'Der Gegner hat richtig erkannt, dass mein erster Vorschlag das Budget deutlich überschreitet. Sein Gegenangebot enthält 28 Gegenstände mit einem Gesamtaufwand von 4+10+50+36+22+68+57+64+56+70+105+167+133+237+224+348+300+434+448+272+382+524+527+528+612+551+689 = 6938 (gerade am Limit). Diese Auswahl hat einen Wert von 103+109+149+135+121+167+156+163+155+169+204+266+232+336+323+447+399+533+547+371+481+623+626+627+711+650+788 = 10963 für mich. Dies ist ein guter Ausgangspunkt, aber ich könnte versuchen, einige niedrigwertige Gegenstände gegen höherwertige mit besserem Verhältnis auszutauschen.'}\\ \tt \\ \tt ARGUMENT: {'Ihr Vorschlag ist budgetkonform und enthält viele wertvolle Gegenstände. Allerdings könnten wir die Gesamtwertigkeit optimieren, indem wir einige niedrigere Wertgegenstände durch höherwertige ersetzen, die ebenfalls effizient sind.'}\\ \tt \\ \tt VORSCHLAG: {'A59', 'A34', 'C12', 'B17', 'A50', 'A99', 'A84', 'C77', 'B07', 'C61', 'B27', 'A23', 'A48', 'C60', 'A42', 'C33', 'C99', 'B49', 'A69', 'C22', 'C23', 'A90', 'B85', 'A62', 'C98', 'C90', 'A87', 'B36', 'C72', 'B54'} 
	  } 
	   } 
	   } 
	 & & \\ 
 

    \theutterance \stepcounter{utterance}  

    & & \multicolumn{4}{p{0.6\linewidth}}{\cellcolor[rgb]{0.9,0.9,0.9}{%
	\makecell[{{p{\linewidth}}}]{% 
	  \tt {\tiny [GM$\rangle$B]}  
	 ARGUMENT: {'Ihr Vorschlag ist budgetkonform und enthält viele wertvolle Gegenstände. Allerdings könnten wir die Gesamtwertigkeit optimieren, indem wir einige niedrigere Wertgegenstände durch höherwertige ersetzen, die ebenfalls effizient sind.'}\\ \tt \\ \tt VORSCHLAG: {'A59', 'A34', 'C12', 'B17', 'A50', 'A99', 'A84', 'C77', 'B07', 'C61', 'B27', 'A23', 'A48', 'C60', 'A42', 'C33', 'C99', 'B49', 'A69', 'C22', 'C23', 'A90', 'B85', 'A62', 'C98', 'C90', 'A87', 'B36', 'C72', 'B54'} 
	  } 
	   } 
	   } 
	 & \\ 
 

    \theutterance \stepcounter{utterance}  

    & & & \multicolumn{4}{p{0.6\linewidth}}{\cellcolor[rgb]{1,0.85,0.72}{%
	\makecell[{{p{\linewidth}}}]{% 
	  \tt {\tiny [GM$\langle$B]}  
	 STRATEGISCHE ÜBERLEGUNGEN: {'Der Gegenspieler hat meinen Vorschlag angenommen und einen modifizierten Vorschlag gemacht, der mehr hochwertige Gegenstände hinzufügt (C90, A87, B36, C72, B54) aber B44 entfernt. Sein Vorschlag hat einen Gesamtaufwand von 4+10+50+36+22+57+64+56+70+105+133+237+224+348+300+434+448+272+382+524+527+528+612+551+689+733+750+804+950+997 = 10463, was immer noch über dem Budget von 6938 liegt. Ich muss dies klarstellen und einen budgetkonformen Gegenvorschlag machen, der meine Top-Prioritäten berücksichtigt. Ich werde die Gegenstände mit dem besten Wert/Aufwand-Verhältnis auswählen.'}\\ \tt \\ \tt ARGUMENT: {'Ihr überarbeiteter Vorschlag überschreitet das Budget immer noch erheblich (10463 > 6938). Lassen Sie uns einen wirklich budgetkonformen Vorschlag mit den effizientesten Gegenständen machen.'}\\ \tt \\ \tt VORSCHLAG: {'A59', 'A34', 'C12', 'B17', 'A50', 'B44', 'A99', 'A84', 'C77', 'B07', 'C61', 'C51', 'B27', 'A23', 'A48', 'C33', 'C99', 'B49', 'A69'} 
	  } 
	   } 
	   } 
	  \\ 
 

    \theutterance \stepcounter{utterance}  

    & & \multicolumn{4}{p{0.6\linewidth}}{\cellcolor[rgb]{0.9,0.9,0.9}{%
	\makecell[{{p{\linewidth}}}]{% 
	  \tt {\tiny [A$\langle$GM]}  
	 ARGUMENT: {'Ihr überarbeiteter Vorschlag überschreitet das Budget immer noch erheblich (10463 > 6938). Lassen Sie uns einen wirklich budgetkonformen Vorschlag mit den effizientesten Gegenständen machen.'}\\ \tt \\ \tt VORSCHLAG: {'A59', 'A34', 'C12', 'B17', 'A50', 'B44', 'A99', 'A84', 'C77', 'B07', 'C61', 'C51', 'B27', 'A23', 'A48', 'C33', 'C99', 'B49', 'A69'} 
	  } 
	   } 
	   } 
	 & \\ 
 

    \theutterance \stepcounter{utterance}  

    & \multicolumn{4}{p{0.6\linewidth}}{\cellcolor[rgb]{0.8,1,0.9}{%
	\makecell[{{p{\linewidth}}}]{% 
	  \tt {\tiny [A$\rangle$GM]}  
	 STRATEGISCHE ÜBERLEGUNGEN: {'Der Gegner hat einen Vorschlag mit 19 Gegenständen gemacht: A59(4), A34(10), C12(50), B17(36), A50(22), B44(68), A99(57), A84(64), C77(56), B07(70), C61(105), C51(167), B27(133), A23(237), A48(224), C33(434), C99(448), B49(272), A69(382). Gesamtaufwand = 4+10+50+36+22+68+57+64+56+70+105+167+133+237+224+434+448+272+382 = 2938. Dies lässt noch 4000 Aufwandseinheiten übrig! Ich kann viele wertvolle Gegenstände hinzufügen. Der Wert für mich ist: 103+109+149+135+121+167+156+163+155+169+204+266+232+336+323+533+547+371+481 = 5220. Ich sollte die effizientesten verbleibenden Gegenstände hinzufügen: C22(524/623=1.19), C23(527/626=1.19), A90(528/627=1.19), B85(612/711=1.16), A62(551/650=1.18), C98(689/788=1.14), C90(733/832=1.14), A87(750/849=1.13), B36(804/903=1.12), C72(950/1049=1.10), B54(997/1096=1.10), A65(906/1005=1.11), C13(945/1044=1.10), A74(873/972=1.11). Ich werde die mit dem besten Verhältnis zuerst hinzufügen.'}\\ \tt \\ \tt ARGUMENT: {'Ihr Vorschlag ist gut, aber wir haben noch erheblichen Budgetspielraum von 4000 Einheiten. Ich schlage vor, mehrere hochwertige Gegenstände hinzuzufügen, die ein ausgezeichnetes Wert-Aufwand-Verhältnis bieten.'}\\ \tt \\ \tt VORSCHLAG: {'A59', 'A34', 'C12', 'B17', 'A50', 'B44', 'A99', 'A84', 'C77', 'B07', 'C61', 'C51', 'B27', 'A23', 'A48', 'C33', 'C99', 'B49', 'A69', 'C22', 'C23', 'A90', 'B85', 'A62', 'C98', 'C90', 'A87', 'B36', 'C72'} 
	  } 
	   } 
	   } 
	 & & \\ 
 

    \theutterance \stepcounter{utterance}  

    & & \multicolumn{4}{p{0.6\linewidth}}{\cellcolor[rgb]{0.9,0.9,0.9}{%
	\makecell[{{p{\linewidth}}}]{% 
	  \tt {\tiny [GM$\rangle$B]}  
	 ARGUMENT: {'Ihr Vorschlag ist gut, aber wir haben noch erheblichen Budgetspielraum von 4000 Einheiten. Ich schlage vor, mehrere hochwertige Gegenstände hinzuzufügen, die ein ausgezeichnetes Wert-Aufwand-Verhältnis bieten.'}\\ \tt \\ \tt VORSCHLAG: {'A59', 'A34', 'C12', 'B17', 'A50', 'B44', 'A99', 'A84', 'C77', 'B07', 'C61', 'C51', 'B27', 'A23', 'A48', 'C33', 'C99', 'B49', 'A69', 'C22', 'C23', 'A90', 'B85', 'A62', 'C98', 'C90', 'A87', 'B36', 'C72'} 
	  } 
	   } 
	   } 
	 & \\ 
 

    \theutterance \stepcounter{utterance}  

    & & & \multicolumn{4}{p{0.6\linewidth}}{\cellcolor[rgb]{1,0.85,0.72}{%
	\makecell[{{p{\linewidth}}}]{% 
	  \tt {\tiny [GM$\langle$B]}  
	 STRATEGISCHE ÜBERLEGUNGEN: {'Der Gegenspieler fügt wieder viele teure Gegenstände hinzu. Ich muss den Gesamtaufwand berechnen: A59(4)+A34(10)+C12(50)+B17(36)+A50(22)+B44(68)+A99(57)+A84(64)+C77(56)+B07(70)+C61(105)+C51(167)+B27(133)+A23(237)+A48(224)+C33(434)+C99(448)+B49(272)+A69(382)+C22(524)+C23(527)+A90(528)+B85(612)+A62(551)+C98(689)+C90(733)+A87(750)+B36(804)+C72(950) = 10463. Das ist immer noch weit über dem Budget von 6938. Ich muss klarstellen, dass wir unter 6938 bleiben müssen, und einen realistischen Vorschlag machen.'}\\ \tt \\ \tt ARGUMENT: {'Ihr Vorschlag überschreitet das Budget von 6938 deutlich (10463). Wir müssen innerhalb des Limits bleiben. Ich schlage einen optimierten Vorschlag vor, der wertvolle Gegenstände priorisiert.'}\\ \tt \\ \tt VORSCHLAG: {'A59', 'A34', 'C12', 'B17', 'A50', 'B44', 'A99', 'A84', 'C77', 'B07', 'C61', 'C51', 'B27', 'A23', 'A48', 'C60', 'C33', 'C99', 'B49', 'A69', 'C22', 'C23', 'A90'} 
	  } 
	   } 
	   } 
	  \\ 
 

    \theutterance \stepcounter{utterance}  

    & & \multicolumn{4}{p{0.6\linewidth}}{\cellcolor[rgb]{0.9,0.9,0.9}{%
	\makecell[{{p{\linewidth}}}]{% 
	  \tt {\tiny [A$\langle$GM]}  
	 ARGUMENT: {'Ihr Vorschlag überschreitet das Budget von 6938 deutlich (10463). Wir müssen innerhalb des Limits bleiben. Ich schlage einen optimierten Vorschlag vor, der wertvolle Gegenstände priorisiert.'}\\ \tt \\ \tt VORSCHLAG: {'A59', 'A34', 'C12', 'B17', 'A50', 'B44', 'A99', 'A84', 'C77', 'B07', 'C61', 'C51', 'B27', 'A23', 'A48', 'C60', 'C33', 'C99', 'B49', 'A69', 'C22', 'C23', 'A90'} 
	  } 
	   } 
	   } 
	 & \\ 
 

    \theutterance \stepcounter{utterance}  

    & \multicolumn{4}{p{0.6\linewidth}}{\cellcolor[rgb]{0.8,1,0.9}{%
	\makecell[{{p{\linewidth}}}]{% 
	  \tt {\tiny [A$\rangle$GM]}  
	 STRATEGISCHE ÜBERLEGUNGEN: {'Der Gegner hat einen Vorschlag mit 23 Gegenständen gemacht. Ich berechne den Aufwand: A59(4) + A34(10) + C12(50) + B17(36) + A50(22) + B44(68) + A99(57) + A84(64) + C77(56) + B07(70) + C61(105) + C51(167) + B27(133) + A23(237) + A48(224) + C60(348) + C33(434) + C99(448) + B49(272) + A69(382) + C22(524) + C23(527) + A90(528) = 4+10+50+36+22+68+57+64+56+70+105+167+133+237+224+348+434+448+272+382+524+527+528 = 5726. Das Budget von 6938 wird nicht überschritten, es bleiben 1212 Einheiten übrig. Der Wert für mich ist: 103+109+149+135+121+167+156+163+155+169+204+266+232+336+323+447+533+547+371+481+623+626+627 = 7223. Ich könnte noch weitere wertvolle Gegenstände hinzufügen, aber das ist bereits ein guter Deal. Ich sollte zustimmen, da dieser Vorschlag budgetkonform ist und mir einen hohen Wert bringt.'}\\ \tt \\ \tt ARGUMENT: {'Dieser Vorschlag ist budgetkonform und enthält viele wertvolle Gegenstände mit gutem Wert-Aufwand-Verhältnis. Er maximiert unseren gemeinsamen Nutzen bei Einhaltung der Budgetbeschränkung.'}\\ \tt \\ \tt ZUSTIMMUNG: {'A59', 'A34', 'C12', 'B17', 'A50', 'B44', 'A99', 'A84', 'C77', 'B07', 'C61', 'C51', 'B27', 'A23', 'A48', 'C60', 'C33', 'C99', 'B49', 'A69', 'C22', 'C23', 'A90'} 
	  } 
	   } 
	   } 
	 & & \\ 
 

    \theutterance \stepcounter{utterance}  

    & & & \multicolumn{2}{p{0.3\linewidth}}{\cellcolor[rgb]{0.95,0.95,0.95}{%
	\makecell[{{p{\linewidth}}}]{% 
	  \tt {\tiny [GM$|$GM]}  
	 game successful 
	  } 
	   } 
	   } 
	 & & \\ 
 

    \theutterance \stepcounter{utterance}  

    & & & \multicolumn{2}{p{0.3\linewidth}}{\cellcolor[rgb]{0.95,0.95,0.95}{%
	\makecell[{{p{\linewidth}}}]{% 
	  \tt {\tiny [GM$|$GM]}  
	 end game 
	  } 
	   } 
	   } 
	 & & \\ 
 

\end{supertabular}
}

\end{document}
